
\subsection{A.5. The universal lift in \S21}

\statement{A.13}{Proved}{The three lift extensions and the mixed fundamental group}
Assume hyperbolic type and choose the additional base point $b$. Write $A=A^+_{g,n}$, $A_0$ for its identity component, $M=\pi_0A$, and $M_*=\pi_0(A^+_{g,n+1,*})$. The topological pullback
\[
 \widetilde A=A\times_M M_*
\]
has the three exact sequences
\[
\begin{split}
1&\longrightarrow\Pi_{g,n}\longrightarrow\widetilde A
                 \longrightarrow A\longrightarrow1,\\
1&\longrightarrow A_0\longrightarrow\widetilde A
                 \longrightarrow M_*\longrightarrow1,\\
1&\longrightarrow A_0\times\Pi_{g,n}\longrightarrow\widetilde A
                 \longrightarrow M\longrightarrow1.
\end{split}                                                   \tag{A.1}
\]
The first extension splits canonically over the subgroup $A_*$ fixing $b$, by lifting maps to fix a chosen point of the universal cover. If $E=EA$ is a universal principal $A$-space (for example the Milnor construction), the associated bundle
$\mathcal X=(E\times U_{g,n})/A_0$ has base $E/A_0$ contractible up to homotopy and fibre $U_{g,n}$. Its residual $M$-action has mixed fundamental group
\[
 \pi_1(\mathcal X\mathbin{/\mkern-6mu/}M)\simeq M_*,
 \qquad
 1\to\pi_1(U_{g,n})\to M_*\to M\to1.                         \tag{A.2}
\]
Here the double slash is the homotopy quotient. These identifications use compatible base points; changing them gives inner conjugation. The same pullback construction applies to the extended groups.
\Man{21}{126--132}{251--263}, especially (74)--(75), scans 252--254.
\begin{proof}
A.3 gives $1\to\Pi\to M_*\to M\to1$. Pull it back along $A\to M$. The first two kernels in (A.1) follow immediately. Over the identity of $M$ the pullback is the \emph{direct} product $A_0\times\Pi$: its multiplication is componentwise. This proves the third sequence and supports the reading $\times$, rather than a nontrivial semidirect product, in (75). Every $f\in A_*$ has a unique lift fixing the chosen lift of $b$; composition of such lifts has the same property. This gives the canonical splitting over $A_*$.

Hyperbolicity implies contractibility of $A_0$, as in the proof of the Birman theorem. Thus $E/A_0$ has the homotopy type of $BA_0$ and is contractible. The associated bundle has fibre $U$, so its fibre inclusion is a homotopy equivalence. Quotienting the residual action gives
$\mathcal X/\!/M\simeq E\times_AU$. Evaluation identifies the homogeneous space $U$ with $A/A_*$; local sections are supplied by isotopy extension. Therefore $E\times_AU\simeq E/A_*\simeq BA_*$. Its fundamental group is $\pi_0A_*=M_*$. The fibration over $BM$ gives exactly (A.2), including the point-pushing kernel. This is the mixed-group construction sketched in \S21, not a genuine action of $M$ on a single marked surface chosen without additional data.
\end{proof}

